% SPDX-License-Identifier: AGPL-3.0-or-later
% Preamble for content/docs/web-first.md: the companion package with boxed
% theorems (so this page's own theorems and proofs show the style), amsthm,
% and the kind the page builds as a tutorial.
\usepackage{amsthm}
\newtheorem{theorem}{Theorem}
\newtheorem{lemma}[theorem]{Lemma}
\newtheorem*{claim}{Claim}
\usepackage[boxed]{reflowtex}

% Like a sectioning command, it leaves the paragraph after it unindented.
\makeatletter
\NewDocumentCommand\pagetitle{o m}{%
  \noindent\IfValueT{#1}{{\small\scshape #1}\par\smallskip\noindent}%
  {\LARGE\bfseries #2\par}\@afterindentfalse\@afterheading}
\makeatother
\newcommand\cs[1]{\texttt{\textbackslash #1}}

% Tutorial: a kind of our own, with a look and no behaviour.
\newenvironment{webwarning}
  {\begin{webstream}{warning}}
  {\end{webstream}}

% Boxed theorems, docs: a boxed copy with its own look.
\DeclareWebBox[keytheorem]{theorem}[accent=#c2410c, class=key]
% Section headings ragged right: at a narrow measure a heading breaks onto
% several lines, and the standard classes would justify them. Only the
% heading's style changes, not its spacing or font.
\usepackage{etoolbox}
\makeatletter
\patchcmd{\@sect}{\begingroup #6}{\begingroup #6\raggedright}{}{}
\patchcmd{\@ssect}{\begingroup #4}{\begingroup #4\raggedright}{}{}
\makeatother
% Character protrusion and font expansion, as in the PDF: the viewer
% protrudes and expands as the document says, and only then.
\usepackage{microtype}
